\chapter{Constant Specht bundles and covolumes}

Let $S=(K,\Omega,\mathcal A,\nu,(|\cdot|_\omega)_\omega)$ be a proper
adelic curve, with $K$ countable of characteristic zero. Write $K_\omega$
for the completed field and $\Omega_\infty$ for the Archimedean places.
The Archimedean absolute values are the usual real or complex absolute
values. Set $m_\infty=\nu(\Omega_\infty)<\infty$.
All tensor norms below are Hilbert norms at infinity and injective norms
at the other places. For a positive-rank adelic bundle, degree is the
negative integral of the logarithmic norm of a nonzero rational determinant
section, and slope is degree divided by rank.

Fix positive integers $a,k$, put $d=ak$, and let $\lambda=(k^a)$ be the
rectangle with $a$ rows and $k$ columns. A word is a function
$w:\{1,\ldots,d\}\to\{1,\ldots,a\}$; its associated basis vector in
$(\Z^a)^{\otimes d}$ is denoted $e_w$. These word vectors are orthonormal
over $\R$. A tableau $t$ is a bijective placement of $1,\ldots,d$ in
$\lambda$. Let $w_t$ record the row of each entry, and let $C_t$ be the
group of permutations preserving each column of $t$. Define
\[
 v_t=\sum_{\sigma\in C_t}\operatorname{sgn}(\sigma)e_{\sigma w_t},
 \qquad L_{a,d}=\sum_{g\in\mathfrak S_d}\Z\,gv_t.
\]
Here permutations act on positions. Relabelling $t$ conjugates its column
group and sends $v_t$ to the correspondingly relabelled vector, so the
lattice is independent of $t$. It is a finite free subgroup of the integral
word space. It has positive rank: the words in the displayed sum are
pairwise distinct, and the coefficient of $e_{w_t}$ is one. Write
$f_{a,d}=\operatorname{rank}_{\Z}L_{a,d}$ and
\[
 c_a(d)=\frac{\log\covol(L_{a,d})}{f_{a,d}}.
\]
The covolume is taken in the real span with the inherited word inner
product. This is the normalization in this paper, Proposition 5.3 and
Lemma 6.1 \cite{bost2026}.

\section{The constant metric and its degree}

\begin{definition}[Constant Specht bundle]
\label{def:constant-specht}
\source{bost-adelic-2026:page-0008,bost-adelic-2026:page-0025}
\lean{BostConjecture.Specht.latticeBundle}
\leanok
Choose an integral basis $b$ of $L_{a,d}$, with vectors $b_1,\ldots,b_f$.
Put $L^S_{a,d}=K\otimes_\Z L_{a,d}$. On its completed fibre define
\[
 \left\|\sum_i z_i b_i\right\|_\omega=
 \begin{cases}
 \displaystyle\left(\sum_w\left|\sum_i z_i b_i(w)\right|_\omega^2\right)^{1/2},
       &\omega\in\Omega_\infty,\\[3pt]
 \displaystyle\max_i|z_i|_\omega,&\omega\notin\Omega_\infty.
 \end{cases}
\]
The resulting constant Specht bundle is denoted $\overline L^S_{a,d}$.
In particular, the matrix at infinity is
$G_b=(\sum_w b_i(w)b_j(w))_{ij}$.
\end{definition}

\begin{lemma}[Integral basis independence]
\label{lem:specht-basis-independence}
\uses{def:constant-specht}
\source{bost-adelic-2026:page-0025}
\lean{BostConjecture.Specht.latticeNormFamily_basis_independent}
\leanok
The completed norms in Definition~\ref{def:constant-specht} are independent
of the integral basis.
\end{lemma}
\begin{proof}
At infinity both expressions are the norm of the same word vector.
At a non-Archimedean place every integer has absolute value at most one.
If $U\in\GL_f(\Z)$ is a basis change, the ultrametric inequality gives
$\|Uz\|_{\max}\le\|z\|_{\max}$. Applying the same inequality to the integral
matrix $U^{-1}$ gives the reverse inequality. Thus either basis has the
same maximum norm.
\end{proof}

\begin{lemma}[Adelic regularity of the constant metric]
\label{lem:specht-regularity}
\uses{def:constant-specht}
\source{bost-adelic-2026:page-0025}
\lean{BostConjecture.Specht.latticeBundle_isHermitian,
BostConjecture.Specht.latticeNormFamily_isStronglyDominated,
BostConjecture.Specht.latticeNormFamily_isMeasurable,
BostConjecture.Specht.latticeBundle}
\leanok
The constant norm family and its dual are measurable and dominated.
Its Archimedean fibres are Hermitian, and its other fibres are ultrametric.
\end{lemma}
\begin{proof}
The Gram matrix $G_b$ is a positive definite integral matrix. The inclusion
of the lattice into word space stays injective after extension to each
$K_\omega$: a nonzero maximal minor over $\Q$ stays nonzero in every
characteristic-zero extension. Thus the displayed word expression is a
positive definite Hermitian norm at infinity. The maximum norm is
ultrametric at every other place.

For a fixed rational vector, each word coordinate and each basis coordinate
belongs to $K$. Its norm is consequently a finite combination of measurable
absolute-value functions, with a measurable choice on $\Omega_\infty$.
For a rational functional $\ell$, density of the rational vectors and
continuity give
\[
 \|\ell\|_{\omega,*}=
 \sup_{0\ne x\in L^S_{a,d}}
       \frac{|\ell(x)|_\omega}{\|x\|_\omega}.
\]
The indexing set is countable. Hence the dual norm is measurable as well.

Compare with the family making $b$ orthonormal at infinity and using its
maximum norm elsewhere. This reference family is dominated by the
adelic integrability of $\log|u|_\omega$ for $u\in K^\times$: finite
coordinate maxima are bounded above by sums of positive parts, and below
by any one nonzero coordinate. The same argument applies to the dual
basis. If $\lambda_{\min},\lambda_{\max}>0$ are the smallest and largest
real eigenvalues of $G_b$, the actual/reference norm ratio at infinity
lies between $\sqrt{\lambda_{\min}}$ and
$\sqrt{\lambda_{\max}}$. These bounds hold over both $\R$ and $\C$.
The inverse bounds hold for the dual ratio. Outside $\Omega_\infty$ the
ratios are one. Their logarithmic distances are therefore bounded by
$C\mathbf1_{\Omega_\infty}$ for a fixed $C$, which is integrable.
This proves domination of the family and its dual.
\end{proof}

\begin{lemma}[Degree of an integral orbit lattice]
\label{lem:constant-lattice-degree}
\uses{lem:specht-regularity}
\source{bost-adelic-2026:page-0026}
\lean{BostConjecture.Specht.latticeBundle_degree,
BostConjecture.Specht.latticeBundle_slope}
\leanok
For any integral Specht orbit lattice $L$ with its word metric and any
integral basis, the constant bundle satisfies
\[
 \widehat{\deg}(\overline L^S)=-m_\infty\log\covol(L),\qquad
 \widehat\mu(\overline L^S)=
   -m_\infty\frac{\log\covol(L)}{\dim_K(K\otimes_\Z L)}.
\]
\end{lemma}
\begin{proof}
Let $b_1,\ldots,b_f$ be the basis and
$\eta=b_1\wedge\cdots\wedge b_f\ne0$. At a finite place the basis and its
dual have norm one, so the determinant norm of $\eta$ is one. At infinity
its squared norm is $\det(G_b)$, and
$\sqrt{\det(G_b)}=\covol(L)$. Consequently
\[
 \log\|\eta\|_\omega
 =\mathbf1_{\Omega_\infty}(\omega)\log\covol(L).
\]
Integrate and negate to obtain the degree, then divide by the dimension
to obtain the slope. Properness makes the degree independent of the
chosen nonzero rational determinant section.
\end{proof}

\begin{lemma}[Rectangular Specht slope]
\label{lem:specht-degree}
\uses{lem:constant-lattice-degree}
\source{bost-adelic-2026:page-0026}
\lean{BostConjecture.Specht.latticeBundle_degree,
BostConjecture.Specht.latticeBundle_slope_rectangle}
\leanok
For $a,k\ge1$ and $d=ak$,
\[
 \widehat{\deg}(\overline L^S_{a,d})
      =-m_\infty\log\covol(L_{a,d}),\qquad
 \widehat\mu(\overline L^S_{a,d})=-m_\infty c_a(d).
\]
\end{lemma}
\begin{proof}
An integral basis remains a basis after extension to $K$, so
$\dim_K(K\otimes_\Z L_{a,d})=f_{a,d}$. Substitute this equality and the
definition of $c_a(d)$ into Lemma~\ref{lem:constant-lattice-degree}.
\end{proof}

\section{Permutation actions and determinant tensors}

\begin{lemma}[Absolute irreducibility of the rectangular orbit span]
\label{lem:rectangular-specht-irreducible}
\source{bost-adelic-2026:page-0025,bost-adelic-2026:page-0026}
\lean{BostConjecture.Arithmetic.Specht.integralRepresentation_baseChange_isIrreducible,
BostConjecture.Arithmetic.Specht.integralRepresentation_baseChange_isAbsolutelyIrreducible}
\leanok
Over every characteristic-zero field $F$, the position action on
$F\otimes_\Z L_{a,d}$ is irreducible. In particular, the action over $K$
is absolutely irreducible.
\end{lemma}
\begin{proof}
Let $M_F$ be the permutation module on balanced words, namely words in
which every letter occurs $k$ times. Give it the bilinear form making
these words orthonormal. For a tableau $t$ set
$\kappa_t=\sum_{\sigma\in C_t}\operatorname{sgn}(\sigma)\sigma$.
If a balanced word repeats a letter within a column of $t$, the
transposition of the two corresponding positions fixes the word and
pairs the summands in $\kappa_t e_w$ with opposite signs. Its image is
therefore zero. Otherwise every column contains all $a$ letters, and
there is a unique $\sigma\in C_t$ sending $w_t$ to $w$. In this case
$\kappa_t e_w=\operatorname{sgn}(\sigma)v_t$. The coefficient of $e_w$
in $v_t$ is the same sign. By linearity,
\[
 \kappa_t x=\langle x,v_t\rangle v_t\qquad(x\in M_F).
\]

The restriction of the bilinear form to $\Q\otimes L_{a,d}$ is
nondegenerate: its Gram matrix in a rational basis is positive definite
over $\R$, hence has nonzero rational determinant. That determinant stays
nonzero over every characteristic-zero field $F$.
If $0\ne U\subseteq F\otimes L_{a,d}$ is invariant, choose $0\ne u\in U$.
Since all $v_t$ span the orbit space, nondegeneracy gives a tableau $t$
with $\langle u,v_t\rangle\ne0$. The displayed identity implies $v_t\in U$.
Its permutation orbit spans the whole space, so $U=F\otimes L_{a,d}$.
This proof also applies to an algebraic closure of $K$.
\end{proof}

\begin{lemma}[Specht semistability and its isometric action]
\label{lem:specht-semistability}
\uses{def:constant-specht,lem:specht-regularity,lem:specht-basis-independence,
lem:rectangular-specht-irreducible,lem:group-semistability}
\source{bost-adelic-2026:page-0008,bost-adelic-2026:page-0026}
\lean{BostConjecture.Specht.latticeBundle_isMeanSlopeSemistable,
BostConjecture.Specht.latticeNormFamily_orbitAction,
BostConjecture.Arithmetic.Specht.integralRepresentation_baseChange_isAbsolutelyIrreducible}
\leanok
The position action of $\mathfrak S_d$ on $\overline L^S_{a,d}$ is
absolutely irreducible and preserves every completed fibre norm.
Consequently $\overline L^S_{a,d}$ is semistable.
\end{lemma}
\begin{proof}
Absolute irreducibility is Lemma~\ref{lem:rectangular-specht-irreducible}.
At infinity a position permutation permutes the orthonormal word
coordinates and hence preserves their squared-norm sum. At any other
place it induces an automorphism of the integral orbit lattice. In an
integral basis both its matrix and inverse have integral entries. The
two ultrametric maximum estimates used in
Lemma~\ref{lem:specht-basis-independence} prove that it is an isometry.
The rank is positive, so the finite-group criterion
Lemma~\ref{lem:group-semistability} gives semistability.
\end{proof}

\begin{lemma}[The determinant map over the original field]
\label{lem:rectangular-determinant-map}
\source{bost-adelic-2026:page-0025}
\lean{BostConjecture.Geometry.freeDeterminantMap,
BostConjecture.Geometry.freeDeterminantMap_basis,
BostConjecture.Geometry.freeDeterminantMap_eq,
BostConjecture.Geometry.freeDeterminantMap_injective,
BostConjecture.Geometry.projectiveDeterminantMap_eq_free}
\leanok
Let $A$ be an $a$-dimensional vector space over a characteristic-zero
field $K$, with $a,k>0$ and $d=ak$. There is a basis-independent injection
\[
 \iota_A:(\det A)^{\otimes k}\otimes_K(K\otimes_\Z L_{a,d})
     \longrightarrow A^{\otimes d}.
\]
For a basis $b$ and $\eta_b=b_1\wedge\cdots\wedge b_a$, the map sends
$\eta_b^{\otimes k}\otimes x$ to the tensor obtained by replacing each
letter $i$ in $x$ by $b_i$. The determinant map obtained by gluing the
same construction on free frames agrees with $\iota_A$.
\end{lemma}
\begin{proof}
For a tableau with consecutive columns, $v_t$ is the product of $k$
alternating determinant tensors. Hence $g^{\otimes d}v_t=(\det g)^kv_t$;
position permutations commute with $g^{\otimes d}$, so the same identity
holds on its entire orbit span. Under $b\mapsto bg$, both the determinant
line generator and this substitution map change by $(\det g)^k$. The
rule is thus independent of $b$. The inclusion of $L_{a,d}$ into the
integral word space remains injective after tensoring with $K$: tensor
first with $\Q$ and then extend scalars to $K$, both of which are exact.
Substitution by the basis is an isomorphism from the resulting word
space to $A^{\otimes d}$. After identifying the determinant line power
with $K$ by $\eta_b^{\otimes k}$, our map is the composite of this
isomorphism with the lattice inclusion. It is therefore injective.
Finally a global basis is itself a free frame. The glued map and
$\iota_A$ have the same expression in that frame, so they agree.
No algebraic closure is used in this construction.
\end{proof}

\begin{lemma}[Rectangular determinant realization of invariants]
\label{lem:rectangular-specht-invariants}
\uses{lem:rectangular-determinant-map}
\source{bost-adelic-2026:page-0025,bost-adelic-2026:page-0026}
\lean{BostConjecture.Geometry.mem_range_freeLatticeMap_iff_specialLinear_fixed,
BostConjecture.Geometry.mem_range_projectiveDeterminantMap_iff}
\leanok
Let $F$ be an algebraically closed field of characteristic zero,
$A$ an $a$-dimensional $F$-vector space, and $a,k>0$, $d=ak$.
For every basis $b$ and every $v\in A^{\otimes d}$,
\[
 v\in\operatorname{im}(b^{\otimes d}|_{F\otimes_\Z L_{a,d}})
 \quad\Longleftrightarrow\quad
 g^{\otimes d}v=v\quad\hbox{for every }g\in\SL(A).
\]
The image of the canonical determinant map equals this lattice image;
this last image equality holds over the original field as well.
Thus the invariant calculation needed for geometric semistability may
be made over $\overline K$, while the determinant map remains defined over $K$.
\end{lemma}
\begin{proof}
Determinant covariance shows that the lattice image is fixed by $\SL(A)$.
Conversely let $v$ be fixed. Given $g\in\GL(A)$, choose $c\in F^\times$
with $c^a=\det g$. Then $c^{-1}g\in\SL(A)$, so
$g^{\otimes d}v=c^dv=(\det g)^kv$.

Let $z$ be the rectangular column-alternating tensor in the basis $b$,
and let $\ell$ be its coefficient covector in the word basis, using
the bilinear coefficient pairing. Each of its $k$ columns has $a!$
distinct signed words, hence $\ell(z)=(a!)^k\ne0$.
Transposing the matrix of $g^{\otimes d}$ and applying determinant
covariance to $z$ gives
$\ell(g^{\otimes d}x)=(\det g)^k\ell(x)$.
The rank-one operator $T(x)=\ell(x)v$ therefore commutes with every
$g^{\otimes d}$. The characteristic-zero Schur--Weyl commutant theorem
expresses $T$ as an $F$-linear combination of position permutations.
Consequently $T(z)=(a!)^kv$ lies in the span of the permutation orbit
of $z$, which is exactly the image of $F\otimes_\Z L_{a,d}$.
Dividing by $(a!)^k$ proves the converse. This is the commutant form of
the Schur--Weyl calculation invoked in Lemma 6.1 of \cite{bost2026}.

In any characteristic-zero field, every determinant-power element is
a scalar multiple of $\eta_b^{\otimes k}$. The formula of
Lemma~\ref{lem:rectangular-determinant-map} therefore shows that its
image is contained in the lattice image. Conversely each substituted
lattice vector is the image of $\eta_b^{\otimes k}\otimes x$.
This proves equality of the two images over the original field,
without making an original-field invariant-space descent assertion.
\end{proof}

\section{The sublinear covolume correction}

For a uniformly chosen standard tableau $T$ of shape $(k^a)$, let
$\lambda^{(t)}$ be the shape of the first $t$ entries. Put
$b_i^{(t)}=\lambda_i^{(t)}+a-i$ for $1\le i\le a$. These shifted row
lengths are strictly decreasing. If entry $j$ is added to row $r_j$,
define
\[
 \delta_{i,r_j}(j)=b_i^{(j-1)}-b_{r_j}^{(j-1)},
 \qquad \gamma_T=\exp\!\left(\sum_{j=1}^d H_j(T)\right),
\]
\[
 H_j(T)=\sum_{i<r_j}\log\frac{\delta_{i,r_j}(j)}{\delta_{i,r_j}(j)-1}.
\]
Legality of insertion implies $\delta_{i,r_j}(j)\ge2$.

\begin{lemma}[The seminormal covolume bound]
\label{lem:specht-gram-bound}
\source{bost-adelic-2026:page-0023,bost-adelic-2026:page-0024}
\lean{BostConjecture.Specht.two_mul_rectangleLogCovolume_le_average_log_prefixNormProduct,
BostConjecture.Specht.rectangleLogCovolume_nonneg}
\leanok
For the uniform average over standard tableaux of shape $(k^a)$,
\[
 0\le 2c_a(ak)\le\mathbb E\log\gamma_T.
\]
\end{lemma}
\begin{proof}
The integral Gram determinant of any lattice basis is a positive integer,
so its logarithm is nonnegative. The standard polytabloids form a rational
basis and belong to the full orbit lattice. Their integral span therefore
has finite index in $L_{a,d}$, and its covolume is this index times
$\covol(L_{a,d})$.
The monic orthogonalization of the standard polytabloids has squared
norms $\gamma_T$, by the seminormal branching formula. In this formula
the contribution when a box is added to row $r$ is exactly
$\prod_{i<r}(b_i-b_r)/(b_i-b_r-1)$; multiplying the contributions along
the tableau gives the displayed definition of $\gamma_T$.
Monic triangular orthogonalization has determinant one, so the Gram
determinant of the standard polytabloids is $\prod_T\gamma_T$.
Taking logarithms of the covolume comparison and dividing by the number
of standard tableaux, which is $f_{a,d}$, proves the upper bound.
Only this comparison is needed; the determinant equality used in this
paper can be replaced by the weaker inequality without changing the
lattice or the limit.
\end{proof}

\begin{lemma}[Bulk separation of shifted rows]
\label{lem:specht-bulk-separation}
\source{bost-adelic-2026:page-0023}
\lean{BostConjecture.Specht.rectanglePrefixGapProbability_uniform}
\leanok
Fix $a\ge2$, $\varepsilon>0$ and $L\ge2$. Uniformly over integers
$\varepsilon ak\le t\le(1-\varepsilon)ak$, the probability that two
shifted row lengths $b_i^{(t)}$ differ by at most $L$ tends to zero as
$k\to\infty$.
\end{lemma}
\begin{proof}
It suffices to consider $0<\varepsilon<1/2$. Set
$c=a(a-1)/2$, $M=k+a-1$, $A=t+c$, $V_0=\prod_{j=0}^{a-1}j!$ and
$V(b)=\prod_{i<j}(b_i-b_j)$. On labelled vectors
$0\le b_i\le M$, $\sum_i b_i=A$, let $B$ be the probability law
\[
 B(b)=\frac{\prod_i\binom M{b_i}}{\binom{aM}A}.
\]
The hook formula applied to the prefix and to the rotated rectangular
complement gives the following two finite identities:
\[
 \mathbb P(\lambda^{(t)}=\lambda)
       =\frac{f^\lambda f^{\lambda^c}}{f^{(k^a)}},\qquad
 \mathbb E_B V(b)^2
       =a!V_0\prod_{j=0}^{a-1}(k+j)^j
           \frac{\binom{ak}t}{\binom{ak+2c}{t+c}}.
\]
Indeed $f^\lambda=t!V(b)/\prod_i b_i!$ and
$f^{\lambda^c}=(ak-t)!V(b)/\prod_i(M-b_i)!$; substitution, and summing
prefix counts to $f^{(k^a)}=(ak)!V_0/\prod_j(k+j)!$, give both formulas.
Thus the tableau law is obtained by weighting $B$ by $V(b)^2$ and sorting.
In the bulk, writing the binomial ratio as rising factorials gives
\[
 \mathbb E_B V(b)^2\ge a!V_0(\varepsilon/a)^{2c}k^c.
\]
To see the inequality, the numerator of that ratio is
$\prod_{r=1}^c(t+r)(ak-t+r)$, whose factors are at least
$\varepsilon ak$; every one of the $2c$ denominator factors is at most
$ak+2c=aM\le a^2k$. Also $\prod_j(k+j)^j\ge k^c$.

The law $B$ counts a uniformly selected $A$-element subset of a population
partitioned into $a$ blocks of size $M$. For any one block, expose this
sample in random order. The martingale of the conditional expected final
block count has increments of absolute value at most one: after $r$
draws its conditional mean is the observed count plus the remaining
sample size times the remaining block proportion; the two possible
next conditional means differ by at most one. The exponential
martingale bound for these bounded increments therefore gives
\[
 B\!\left(\max_i|b_i-A/a|\ge u\right)
       \le2a\exp\left(-\frac{u^2}{2a^2k}\right),
\]
since $A\le aM\le a^2k$. If the whole population is sampled the counts
are deterministic and the inequality is immediate.
On the complementary event, all pairwise gaps are less than $2u$.
If some gap is at most $L$, then
$V(b)^2\le L^2(2u)^{2(c-1)}$; everywhere $V(b)^2\le(ak)^{2c}$.
Dividing these two numerator bounds by the displayed lower bound for
$\mathbb E_B V(b)^2$ bounds the desired probability by
\[
 \frac{L^2(2u)^{2(c-1)}
        +2a(ak)^{2c}\exp(-u^2/(2a^2k))}
      {a!V_0(\varepsilon/a)^{2c}k^c}.
\]
Take $u=k^{1/2+1/(4c)}$. The first term is a constant times
$k^{-(c+1)/(2c)}$, and the second is a constant times
$k^c\exp(-k^{1/(2c)}/(2a^2))$. Both tend to zero, independently of $t$.
The small-gap event and these bounds are invariant under sorting.
\end{proof}

\begin{lemma}[Sublinear rectangular covolume]
\label{lem:specht-limit}
\uses{lem:specht-gram-bound,lem:specht-bulk-separation}
\source{bost-adelic-2026:page-0023,bost-adelic-2026:page-0024}
\lean{BostConjecture.Specht.tendsto_rectangleLogCovolume_div}
\leanok
For every fixed integer $a\ge1$,
\[
 \lim_{k\to\infty}\frac{c_a(ak)}{ak}=0.
\]
\end{lemma}
\begin{proof}
If $a=1$, the lattice is generated by its single unit word, so $c_1(k)=0$.
Suppose $a\ge2$, and put $C=(a-1)\log2$. Each legal insertion satisfies
$0\le H_j\le C$. If all the shifted gaps of its prefix exceed $L\ge2$,
then $\log(1+x)\le x$ for $x\ge0$ gives
$H_j\le(a-1)/(L-1)$. This bound holds for every possible next row and
therefore depends only on the prefix, without conditioning on that row.
Let $p_{k,\varepsilon,L}$ be the largest bulk small-gap probability.
There are at most $2\varepsilon ak+2$ nonbulk times. Splitting the
expectation at each bulk time into its small-gap event and complement,
then summing, gives
\[
 0\le\frac{\mathbb E\log\gamma_T}{ak}
 \le\left(2\varepsilon+\frac2{ak}\right)C
       +\frac{a-1}{L-1}+C p_{k,\varepsilon,L}.
\]
Lemma~\ref{lem:specht-bulk-separation} makes the last term tend to zero.
First take the superior limit in $k$, then let $L\to\infty$, and finally
$\varepsilon\downarrow0$. The resulting superior limit is zero.
Lemma~\ref{lem:specht-gram-bound} and nonnegativity now squeeze
$c_a(ak)/(ak)$ to zero. No effective rate is required.
\end{proof}
